\chapter{KESIMPULAN DAN SARAN}

Berdasarkan pembahasan pada Bab 4, dapat ditarik kesimpulan sebagai berikut.

\begin{enumerate}
      \item Syarat perlu komutativitas dua Latin square terhadap perkalian max-plus
            diperoleh dari elemen maksimum. Berdasarkan
            \autoref{thm:komutatif_elemen_maksimum}, jika $A,B\in L_S(n)$
            komutatif, maka permutasi yang bersesuaian dengan simbol maksimum
            harus saling komutatif. Dengan demikian, kandidat permutasi simbol
            maksimum pada $B$ harus berada dalam centralizer dari permutasi simbol
            maksimum pada $A$.

      \item Kriteria komutativitas dua Latin square terhadap perkalian max-plus
            diperoleh melalui syarat cukup dan kriteria perlu-cukup. Berdasarkan
            \autoref{thm:syarat_cukup_subgrup_abelian}, komutativitas dijamin
            ketika seluruh permutasi penyusun $A$ dan $B$ berada dalam subgrup
            abelian yang sama. Selanjutnya, berdasarkan
            \autoref{thm:kriteria_komutatif_superlevel}, komutativitas terjadi
            jika dan hanya jika himpunan superlevel pada $A\otimes B$ dan
            $B\otimes A$ sama untuk setiap level yang diperlukan.

      \item Konstruksi pasangan Latin square yang komutatif dengan suatu Latin
            square $A$ diberikan pada
            \autoref{alg:pencarian_pasangan_latin_square_komutatif}. Algoritma
            tersebut memanfaatkan tahap alternatif dari
            \autoref{thm:syarat_cukup_subgrup_abelian} serta tahap umum melalui
            centralizer dan pemeriksaan superlevel. Penerapan pada ordo $3$, $4$,
            dan $5$ menunjukkan bahwa algoritma tersebut dapat digunakan untuk
            memperoleh kandidat $B$ yang memenuhi $A\otimes B=B\otimes A$.
\end{enumerate}

Berdasarkan pembahasan sebelumnya, berikut merupakan saran untuk penelitian
selanjutnya.

\begin{enumerate}
      \item Mengimplementasikan
            \autoref{alg:pencarian_pasangan_latin_square_komutatif} secara
            komputasi untuk menguji Latin square berordo lebih besar.

      \item Mengkaji efisiensi pemeriksaan superlevel agar proses pencarian
            kandidat dapat dilakukan lebih cepat.

      \item Mengembangkan syarat cukup melalui subgrup abelian ke bentuk yang
            lebih umum, misalnya melalui centralizer dari subgrup permutasi
            penyusun Latin square.
\end{enumerate}